% !PRINTED 266
\weberpage{266}

%<<CONT
und die Wege kennen gelernt, wie man auch in anderen Fällen zu ihrer Berechnung gelangen kann.

Wir haben schon früher bemerkt, dass die Bezoutiante nicht eine allgemeine quadratische Form ist, sondern dass sie die besondere Eigenthümlichkeit hat, dass die Anzahl $\nu$ ihrer negativen Quadrate niemals grösser als $\frac{1}{2}n$ werden kann. Diese Eigenthümlichkeit muss in gewissen Ungleichheitsbedingungen zwischen den Co\"efficienten ihren Ausdruck finden, die aber zur Zeit noch nicht bekannt sind. Nur in dem Falle $n = 3$ können wir den algebraischen Charakter dieser Beschränkung vollständig angeben.

Für die cubische Gleichung
\[
f(x) = a_0 x^3 + a_1 x^2 + a_2 x + a_3 = 0
\]
ist die Bezoutiante nach \S.\,73, (13) nichts Anderes, als die \gesperrt{Hessesche} Covariante
\[
-\tfrac{2}{3} H(t_1, t_0) = -\tfrac{2}{3}(A_0 t_1^2 + A_1 t_1 t_0 + A_2 t_0^2).
\]
Es besteht aber, wenn $D$ die Discriminante, also
\[
3D = 4A_0 A_2 - A_1^2
\]
ist, die Relation [\S.\,62, (4), (12)]
\[
4H^3 + Q^2 + 27Df^2 = 0,
\]
mithin, wenn wir darin $t_0 = 0$ setzen,

\weq{1}{-4A_0^3 + q_0^2 + 27Da_0^2 = 0,}

wenn
\[
q_0 = 27a_0^2 a_3 - 9a_0 a_1 a_2 + 2a_1^3
\]
ist. Diese Relation zeigt, dass $A_0$ und $D$ nicht zugleich positiv sein können; und dass also die Form $-H$ nicht in zwei negative Quadrate zerlegbar ist.

Wir wollen den allgemeinen Satz des vorigen Paragraphen noch auf die biquadratische Gleichung anwenden, die wir der Einfachheit halber in der Form

\weq{2}{x^4 + ax^2 + bx + c = 0}

annehmen. Um die Co\"efficienten der Bezoutiante zu bilden, haben wir in den Formeln \S.\,73, (14)
\[
a_0,\ a_1,\ a_2,\ a_3,\ a_4
\]
durch
\[
1,\ 0,\ a,\ b,\ c
\]
zu ersetzen. Wir erhalten so
\[
\begin{aligned}
&B_{2,2} = -2a, \quad B_{1,1} = a^2 - 4c, \quad B_{0,0} = -\tfrac{1}{4}(8ac - 3b^2), \\
&B_{0,1} = \tfrac{1}{2}ab, \quad B_{2,0} = -4c, \quad B_{1,2} = -3b.
\end{aligned}
\]
% [resolved 2026-08-06, 8x zoom strip 8/8] The smudged sign inside B_{0,0} is a MINUS.
% Compared against the clean minus signs on these same two lines (-2a, -4c, -3b, -1/4) and
% the plus signs three lines above (-4A_0^3 + q_0^2 + 27 D a_0^2): the disputed stroke has
% no vertical bar, only a horizontal one with a speck of ink. B_{0,0} = -1/4(8ac - 3b^2)
% stands as transcribed. Same strip also confirms B_{2,2} = -2a, B_{1,1} = a^2 - 4c,
% B_{0,1} = 1/2 ab, B_{2,0} = -4c, B_{1,2} = -3b.
